\begin{statementbox}

\textbf{\textcolor{CStatement}{条件}}
\begin{description}[style=nextline, leftmargin=2.8em, labelsep=0.5em, itemsep=0.35em]
  \item[\textbf{\textcolor{CStatement}{条件1（焦点在 $x$ 轴）：}}]
  双曲线 $\displaystyle\frac{x^2}{a^2}-\frac{y^2}{b^2}=1$，其中
  $a>0,b>0$，$c=\sqrt{a^2+b^2}$，$e=\frac ca>1$；
  左焦点 $F_1(-c,0)$，右焦点 $F_2(c,0)$。
  \item[\textbf{\textcolor{CStatement}{条件2（点在曲线上）：}}]
  $P(x_0,y_0)$ 在双曲线上，故 $|x_0|\ge a$：
  右支 $x_0\ge a$，左支 $x_0\le-a$。
\end{description}

\medskip
\textbf{\textcolor{CStatement}{结论}}
\begin{description}[style=nextline, leftmargin=2.8em, labelsep=0.5em, itemsep=0.45em]
  \item[\textbf{\textcolor{CStatement}{结论一（统一焦半径公式）：}}]
  同一点到两个焦点的距离，只用横坐标就能计算：
  \[
    \boxed{|PF_1|=|a+ex_0|},\qquad
    \boxed{|PF_2|=|a-ex_0|}.
  \]

  \item[\textbf{\textcolor{CStatement}{结论二（按左右支去绝对值）：}}]
  \[
    \begin{cases}
      |PF_1|=ex_0+a,\quad |PF_2|=ex_0-a,&x_0\ge a\quad(\text{右支}),\\[.4em]
      |PF_1|=-ex_0-a,\quad |PF_2|=-ex_0+a,&x_0\le-a\quad(\text{左支}).
    \end{cases}
  \]

  \item[\textbf{\textcolor{CStatement}{结论三（更容易理解的对称记法）：}}]
  记 $t=e|x_0|$，则 $t\ge c>a$，
  \[
    \boxed{\text{到同侧焦点的距离}=t-a},\qquad
    \boxed{\text{到异侧焦点的距离}=t+a}.
  \]
  因而两焦半径\textbf{差的绝对值为 $2a$}，\textbf{和为 $2e|x_0|$}。
\end{description}

\medskip
\textbf{\textcolor{CStatement}{使用提醒：}}
\begin{itemize}[leftmargin=*,itemsep=0.25em,topsep=0.3em]
  \item 先辨左、右支，再去绝对值。双曲线上不存在 $-a<x_0<a$ 的点。
  \item $F_1$ 永远是左焦点，$F_2$ 永远是右焦点；\textbf{同侧/异侧}随支改变。
  \item 在顶点 $x_0=\pm a$ 时也成立；右顶点处 $(|PF_1|,|PF_2|)=(c+a,c-a)$，左顶点处反过来。
\end{itemize}

\end{statementbox}
